\subsection{Mutating External Scopes}

Reading a name from an outer scope and assigning to a name in an outer scope are different operations. Python allows ordinary read access through the LEGB chain, but assignment has stricter rules.

By default, assignment inside a function creates or updates a local binding in that function frame. It does not automatically rewrite a global name or an enclosing function's name.

There are two explicit declarations for overriding this default:

\begin{itemize}
    \item \texttt{global} targets the module namespace;
    \item \texttt{nonlocal} targets an enclosing function scope.
\end{itemize}

These declarations do not change the object-passing model. They change which namespace receives the binding operation.

\subsubsection{Local Read Access Boundaries vs. the Shadowing Consequence of Assignment}

A function can read a global name without a declaration.

\begin{verbatim}
answer = 42

def show_answer():
    print(f"answer inside function: {answer}")

show_answer()
print(f"answer outside function: {answer}")
\end{verbatim}

Possible output:

\begin{verbatim}
answer inside function: 42
answer outside function: 42
\end{verbatim}

The function body does not define a local \texttt{answer}, so Python searches outward and finds the global \texttt{answer}.

Assignment is different.

\begin{verbatim}
answer = 42

def set_answer():
    answer = 100
    print(f"inside function: {answer}")

set_answer()
print(f"outside function: {answer}")
\end{verbatim}

Possible output:

\begin{verbatim}
inside function: 100
outside function: 42
\end{verbatim}

The assignment inside \texttt{set\_answer()} creates a local name. It does not mutate the module-level name.

The structure is:

\begin{verbatim}
module namespace:
    answer -> 42

set_answer() local frame:
    answer -> 100
\end{verbatim}

The two names are separate bindings.

This is shadowing: the local name hides the outer name during lookup inside the function.

\begin{verbatim}
message = "global message"

def show_shadowing():
    message = "local message"

    print(f"local lookup: {message}")
    print(f"'message' in locals():  {'message' in locals()}")
    print(f"'message' in globals(): {'message' in globals()}")

show_shadowing()

print(f"global lookup: {message}")
\end{verbatim}

Possible output:

\begin{verbatim}
local lookup: local message
'message' in locals():  True
'message' in globals(): True
global lookup: global message
\end{verbatim}

The name \texttt{message} exists in both namespaces. Inside the function, the local binding wins because Local comes before Global in LEGB.

The local namespace display is a dictionary-like object.

\begin{verbatim}
message = "global message"

def show_shadowing():
    message = "local message"
    local_map = locals()

    print(f"type(local_map): {type(local_map)}")
    print()

    selected_names = [
        "keys",
        "values",
        "items",
        "get",
        "__contains__",
        "__getitem__",
    ]

    for name in selected_names:
        print(f"{name:<14} {name in dir(local_map)}")

    print()
    print(f"local_map: {local_map}")

show_shadowing()
\end{verbatim}

Possible output:

\begin{verbatim}
type(local_map): <class 'dict'>

keys           True
values         True
items          True
get            True
__contains__   True
__getitem__    True

local_map: {'message': 'local message'}
\end{verbatim}

Reading an outer mutable object and mutating that object in place does not require rebinding the outer name.

\begin{verbatim}
quotes = ["D'oh!"]

def add_quote():
    quotes.append("No soup for you!")

add_quote()

print(f"quotes: {quotes}")
\end{verbatim}

Possible output:

\begin{verbatim}
quotes: ["D'oh!", 'No soup for you!']
\end{verbatim}

The function did not assign to the name \texttt{quotes}. It read the global binding, obtained the list object, and called the list object's \texttt{append()} method.

The list object is mutable.

\begin{verbatim}
quotes = ["D'oh!"]

print(f"type(quotes): {type(quotes)}")
print()

selected_names = [
    "append",
    "extend",
    "insert",
    "pop",
    "clear",
    "__len__",
    "__getitem__",
    "__setitem__",
]

for name in selected_names:
    print(f"{name:<14} {name in dir(quotes)}")

print()
quotes.append("The answer is 42.")
print(f"after append: {quotes}")
\end{verbatim}

Possible output:

\begin{verbatim}
type(quotes): <class 'list'>

append         True
extend         True
insert         True
pop            True
clear          True
__len__        True
__getitem__    True
__setitem__    True

after append: ["D'oh!", 'The answer is 42.']
\end{verbatim}

This operation changes the shared list object. It does not reassign the global name.

This does reassign the local name:

\begin{verbatim}
quotes = ["D'oh!"]

def replace_quotes():
    quotes = ["This is an ex-list."]
    print(f"inside function: {quotes}")

replace_quotes()

print(f"outside function: {quotes}")
\end{verbatim}

Possible output:

\begin{verbatim}
inside function: ['This is an ex-list.']
outside function: ["D'oh!"]
\end{verbatim}

So the basic distinction is:

\begin{verbatim}
quotes.append(x)
    read outer name
    mutate shared list object

quotes = [...]
    bind local name unless global or nonlocal says otherwise
\end{verbatim}

This distinction explains why some functions appear to ``change globals'' without \texttt{global}. They are not changing the global binding. They are mutating an object reached through that binding.

\subsubsection{Overriding Module Scope: The \texttt{global} Declaration Syntax and Module Namespace Mutation}

The \texttt{global} declaration tells Python that assignments to a name inside the function should target the module namespace.

\begin{verbatim}
answer = 42

def set_global_answer():
    global answer

    answer = 100
    print(f"inside function: {answer}")

print(f"before call: {answer}")

set_global_answer()

print(f"after call:  {answer}")
\end{verbatim}

Possible output:

\begin{verbatim}
before call: 42
inside function: 100
after call:  100
\end{verbatim}

The assignment:

\begin{verbatim}
answer = 100
\end{verbatim}

does not create a local binding here. Because of \texttt{global answer}, it updates the module namespace binding.

The module namespace can be inspected with \texttt{globals()}.

\begin{verbatim}
answer = 42

def set_global_answer():
    global answer

    print(f"before assignment in globals: {globals()['answer']}")
    answer = 100
    print(f"after assignment in globals:  {globals()['answer']}")

set_global_answer()
\end{verbatim}

Possible output:

\begin{verbatim}
before assignment in globals: 42
after assignment in globals:  100
\end{verbatim}

The object returned by \texttt{globals()} is the actual module namespace dictionary.

\begin{verbatim}
answer = 42

global_map = globals()

print(f"type(global_map): {type(global_map)}")
print()

selected_names = [
    "keys",
    "values",
    "items",
    "get",
    "__contains__",
    "__getitem__",
    "__setitem__",
]

for name in selected_names:
    print(f"{name:<14} {name in dir(global_map)}")

print()
print(f"global_map['answer']: {global_map['answer']}")
\end{verbatim}

Possible output:

\begin{verbatim}
type(global_map): <class 'dict'>

keys           True
values         True
items          True
get            True
__contains__   True
__getitem__    True
__setitem__    True

global_map['answer']: 42
\end{verbatim}

The \texttt{global} declaration is needed only for rebinding the name. It is not needed for mutating an already reachable object.

\begin{verbatim}
quotes = ["D'oh!"]

def mutate_global_list():
    quotes.append("No soup for you!")

def rebind_global_list():
    global quotes
    quotes = ["It is only a model."]

print(f"start: {quotes}")

mutate_global_list()
print(f"after mutation: {quotes}")

rebind_global_list()
print(f"after rebinding: {quotes}")
\end{verbatim}

Possible output:

\begin{verbatim}
start: ["D'oh!"]
after mutation: ["D'oh!", 'No soup for you!']
after rebinding: ['It is only a model.']
\end{verbatim}

The first function changes the list object. The second function changes the module-level binding \texttt{quotes} itself.

The \texttt{global} declaration must appear before the name is used in a way that conflicts with the declaration inside the same function body.

This is valid:

\begin{verbatim}
score = 42

def set_score():
    global score
    score = 100

set_score()
print(score)
\end{verbatim}

Possible output:

\begin{verbatim}
100
\end{verbatim}

A \texttt{global} declaration should be used sparingly. It creates hidden coupling between the function body and the module state. When possible, it is usually clearer to return a value and let the caller decide whether to rebind the global name.

\begin{verbatim}
score = 42

def make_new_score(old_score):
    return old_score + 58

score = make_new_score(score)

print(score)
\end{verbatim}

Possible output:

\begin{verbatim}
100
\end{verbatim}

This version makes the state transfer explicit:

\begin{verbatim}
caller passes old value
function returns new value
caller rebinds module name
\end{verbatim}

The \texttt{global} version directly mutates the module namespace from inside the function.

\subsubsection{Overriding Nested Intermediary Scopes: The \texttt{nonlocal} Declaration Syntax and Cell Reference Mutation}

The \texttt{nonlocal} declaration targets a name in an enclosing function scope. It does not target the module namespace.

Without \texttt{nonlocal}, assignment inside an inner function creates a new local binding.

\begin{verbatim}
def outer():
    count = 0

    def inner():
        count = 1
        print(f"inner count: {count}")

    inner()
    print(f"outer count: {count}")

outer()
\end{verbatim}

Possible output:

\begin{verbatim}
inner count: 1
outer count: 0
\end{verbatim}

The \texttt{count = 1} assignment creates a local \texttt{count} inside \texttt{inner()}. It does not modify \texttt{outer()}'s \texttt{count}.

With \texttt{nonlocal}, the assignment targets the nearest enclosing function scope that already has that name.

\begin{verbatim}
def outer():
    count = 0

    def inner():
        nonlocal count

        count = count + 1
        print(f"inner count: {count}")

    inner()
    inner()
    inner()

    print(f"outer count: {count}")

outer()
\end{verbatim}

Possible output:

\begin{verbatim}
inner count: 1
inner count: 2
inner count: 3
outer count: 3
\end{verbatim}

The inner function modifies the binding preserved in the enclosing scope.

This is especially important for closures.

\begin{verbatim}
def make_counter():
    count = 0

    def next_count():
        nonlocal count

        count = count + 1
        return count

    return next_count

counter = make_counter()

print(counter())
print(counter())
print(counter())
\end{verbatim}

Possible output:

\begin{verbatim}
1
2
3
\end{verbatim}

The outer function has already returned, but the inner function still carries access to \texttt{count} through a closure cell.

The closure can be inspected.

\begin{verbatim}
def make_counter():
    count = 0

    def next_count():
        nonlocal count

        count = count + 1
        return count

    return next_count

counter = make_counter()

print(f"counter.__closure__: {counter.__closure__}")
print(f"counter.__code__.co_freevars: {counter.__code__.co_freevars}")
print()

cell = counter.__closure__[0]

print(f"cell:       {cell}")
print(f"type(cell): {type(cell)}")
print(f"contents:   {cell.cell_contents}")
\end{verbatim}

Possible output:

\begin{verbatim}
counter.__closure__: (<cell at 0x...: int object at 0x...>,)
counter.__code__.co_freevars: ('count',)

cell:       <cell at 0x...: int object at 0x...>
type(cell): <class 'cell'>
contents:   0
\end{verbatim}

After calls, the cell contents change.

\begin{verbatim}
def make_counter():
    count = 0

    def next_count():
        nonlocal count

        count = count + 1
        return count

    return next_count

counter = make_counter()
cell = counter.__closure__[0]

print(f"before: {cell.cell_contents}")

print(counter())
print(counter())

print(f"after:  {cell.cell_contents}")
\end{verbatim}

Possible output:

\begin{verbatim}
before: 0
1
2
after:  2
\end{verbatim}

The cell object can be inspected.

\begin{verbatim}
def make_counter():
    count = 0

    def next_count():
        nonlocal count

        count = count + 1
        return count

    return next_count

counter = make_counter()
cell = counter.__closure__[0]

print(f"type(cell): {type(cell)}")
print()

selected_names = [
    "cell_contents",
    "__class__",
    "__eq__",
    "__repr__",
]

for name in selected_names:
    print(f"{name:<16} {name in dir(cell)}")

print()
print(f"cell.cell_contents: {cell.cell_contents}")
\end{verbatim}

Possible output:

\begin{verbatim}
type(cell): <class 'cell'>

cell_contents    True
__class__        True
__eq__           True
__repr__         True

cell.cell_contents: 0
\end{verbatim}

The \texttt{nonlocal} declaration changes which binding is updated. It does not mean that integers become mutable. In the counter example, this line:

\begin{verbatim}
count = count + 1
\end{verbatim}

creates or obtains a new integer object, then updates the enclosing cell so that \texttt{count} refers to the new object.

A mutable object in an enclosing scope can be mutated without \texttt{nonlocal}, just like a global mutable object can be mutated without \texttt{global}.

\begin{verbatim}
def outer():
    lines = []

    def inner():
        lines.append("Ni!")

    inner()
    inner()

    print(lines)

outer()
\end{verbatim}

Possible output:

\begin{verbatim}
['Ni!', 'Ni!']
\end{verbatim}

The inner function reads the enclosing binding \texttt{lines}, obtains the list object, and mutates the list.

But rebinding the enclosing name requires \texttt{nonlocal}.

\begin{verbatim}
def outer():
    lines = []

    def inner():
        nonlocal lines

        lines = ["This is an ex-list."]

    inner()

    print(lines)

outer()
\end{verbatim}

Possible output:

\begin{verbatim}
['This is an ex-list.']
\end{verbatim}

The difference is the same as before:

\begin{verbatim}
lines.append("Ni!")
    mutate the shared list object

lines = [...]
    rebind the enclosing name because nonlocal was declared
\end{verbatim}

The \texttt{nonlocal} name must already exist in an enclosing function scope. It cannot create a new enclosing name, and it does not refer to the global namespace.

A useful practical rule is:

\begin{quote}
Use \texttt{nonlocal} only when an inner function must rebind a name from an enclosing function. For ordinary data flow, prefer returning values unless closure state is intentionally part of the design.
\end{quote}

\subsubsection{\texttt{locals()} Caveats: Why the Displayed Local Namespace Is Not Always a Directly Writable Control Surface}

The function \texttt{locals()} returns a mapping that represents the current local namespace. At module level, local and global namespace are the same namespace.

\begin{verbatim}
answer = 42

print(f"locals() is globals(): {locals() is globals()}")

locals()["answer"] = 100

print(f"answer: {answer}")
\end{verbatim}

Possible output at module level:

\begin{verbatim}
locals() is globals(): True
answer: 100
\end{verbatim}

At module level, \texttt{locals()} and \texttt{globals()} usually refer to the same dictionary. Mutating that dictionary mutates the module namespace.

Inside an ordinary function, this should not be treated as a reliable way to create or rebind local variables.

\begin{verbatim}
def try_write_locals():
    x = 42

    local_map = locals()
    local_map["x"] = 100
    local_map["y"] = "D'oh!"

    print(f"local_map: {local_map}")
    print(f"x:         {x}")

    try:
        print(f"y:         {y}")
    except NameError as err:
        print(f"error type: {type(err)}")
        print(f"message:    {err}")

try_write_locals()
\end{verbatim}

Possible output:

\begin{verbatim}
local_map: {'x': 100, 'y': "D'oh!"}
x:         42
error type: <class 'NameError'>
message:    name 'y' is not defined
\end{verbatim}

The displayed dictionary changed, but the actual fast-local variable \texttt{x} used by the function body did not necessarily change. The name \texttt{y} did not become a normal local variable.

This is the caveat: inside optimized function frames, \texttt{locals()} is mainly an inspection surface, not a dependable write-control surface.

The map returned by \texttt{locals()} may look writable because it is dictionary-like.

\begin{verbatim}
def inspect_locals_map():
    x = 42
    local_map = locals()

    print(f"type(local_map): {type(local_map)}")
    print()

    selected_names = [
        "keys",
        "values",
        "items",
        "get",
        "__getitem__",
        "__setitem__",
    ]

    for name in selected_names:
        print(f"{name:<14} {name in dir(local_map)}")

inspect_locals_map()
\end{verbatim}

Possible output:

\begin{verbatim}
type(local_map): <class 'dict'>

keys           True
values         True
items          True
get            True
__getitem__    True
__setitem__    True
\end{verbatim}

But dictionary-like appearance does not guarantee that assigning into it rewires the optimized local slots of the running function.

The reason connects to the code object's local-variable layout.

\begin{verbatim}
def sample_func():
    x = 42
    y = "D'oh!"
    return x, y

code_obj = sample_func.__code__

print(f"co_varnames: {code_obj.co_varnames}")
print(f"co_nlocals:  {code_obj.co_nlocals}")
\end{verbatim}

Possible output:

\begin{verbatim}
co_varnames: ('x', 'y')
co_nlocals:  2
\end{verbatim}

CPython knows the ordinary local variable slots ahead of time. The actual execution of \texttt{x} and \texttt{y} uses that optimized local layout, not necessarily a live dictionary lookup for every access.

This is why writing into \texttt{locals()} inside a function is not a normal assignment substitute.

Use ordinary assignment for local names:

\begin{verbatim}
def correct_local_assignment():
    x = 42
    x = 100
    print(f"x: {x}")

correct_local_assignment()
\end{verbatim}

Possible output:

\begin{verbatim}
x: 100
\end{verbatim}

Use \texttt{global} for module-level rebinding:

\begin{verbatim}
answer = 42

def correct_global_assignment():
    global answer

    answer = 100

correct_global_assignment()
print(answer)
\end{verbatim}

Possible output:

\begin{verbatim}
100
\end{verbatim}

Use \texttt{nonlocal} for enclosing-function rebinding:

\begin{verbatim}
def outer():
    answer = 42

    def inner():
        nonlocal answer

        answer = 100

    inner()
    print(answer)

outer()
\end{verbatim}

Possible output:

\begin{verbatim}
100
\end{verbatim}

The rule is:

\begin{quote}
Use \texttt{locals()} to inspect local bindings. Do not use it as the normal mechanism for changing local variables inside an optimized function frame.
\end{quote}

The scope-mutation summary is:

\begin{center}
\begin{tabular}{|p{0.30\linewidth}|p{0.32\linewidth}|p{0.28\linewidth}|}
\hline
\textbf{Goal} & \textbf{Correct mechanism} & \textbf{Typical effect} \\
\hline
Create or update a local name & ordinary assignment & updates current function frame \\
\hline
Rebind a module-level name & \texttt{global} declaration & updates module namespace dictionary \\
\hline
Rebind an enclosing function name & \texttt{nonlocal} declaration & updates closure cell / enclosing binding \\
\hline
Mutate a reachable mutable object & method call or item assignment & changes the shared object, not the name binding \\
\hline
Inspect current locals & \texttt{locals()} & displays local namespace information \\
\hline
\end{tabular}
\end{center}

The central distinction remains unchanged:

\begin{quote}
Name rebinding changes a namespace entry. Object mutation changes an object reached through a name.
\end{quote}